\documentclass[10pt]{article}
\usepackage[a4paper,margin=26mm]{geometry}
\usepackage[T1]{fontenc}
\usepackage{lmodern}
\usepackage{amsmath,amssymb}
\emergencystretch=2em
\title{A purely odd representation of negative superdimension\\Disproof of Conjecture 00000007400}
\author{ziangni-sys}
\date{4 October 2026}
\begin{document}
\maketitle
\section*{Claim and conventions}
The original conjecture asserts nonnegativity of superdimensions of
supersymmetric representations. Neither the English nor Chinese statement
restricts the grading, excludes trivial actions, or imposes a positivity
condition on the character.
For a finite-dimensional complex super vector space
$V=V_{\bar0}\oplus V_{\bar1}$, its standard superdimension is
\[
\operatorname{sdim}V=\dim_{\mathbb C}V_{\bar0}
-\dim_{\mathbb C}V_{\bar1}.
\]
Unlike ordinary dimension, this difference can be negative.

\section*{A complete superrepresentation}
Let $\mathfrak g$ be the one-dimensional abelian complex Lie algebra
viewed as a Lie superalgebra concentrated in even degree:
$\mathfrak g_{\bar0}=\mathbb C$, $\mathfrak g_{\bar1}=0$, and bracket zero.
Bilinearity, graded skew symmetry and the graded Jacobi identity all hold
because every bracket vanishes.

Let the underlying vector space of $V$ be $\mathbb C$, with grading
$V_{\bar0}=\{0\}$ and $V_{\bar1}=\mathbb C$.
These are actual complex subspaces, disjoint and summing to $V$.
Define the action $\rho(a)v=0$ for every $a\in\mathfrak g$ and $v\in V$.
It is bilinear and preserves parity. Its representation identity is
\[
\rho([a,b])v
=\rho(a)\rho(b)v-\rho(b)\rho(a)v=0.
\]
Since all nonzero algebra elements are even, this is exactly the required
superrepresentation identity; requirements involving odd algebra elements
hold trivially.
Thus $V$ is a finite-dimensional graded representation, not merely an
assigned pair of dimension counts.

Its dimensions are $\dim V_{\bar0}=0$ and $\dim V_{\bar1}=1$, so
\[
\operatorname{sdim}V=0-1=-1<0.
\]
This disproves the universal nonnegativity assertion.
The same example is the parity reversal of the one-dimensional trivial
even representation.
Its supercharacter at the identity is the supertrace of the identity,
equal to $-1$; character terminology does not turn it into a positive
number.

\section*{Formal verification}
The Lean project models the actual complex subspaces and their direct
decomposition, the zero bilinear Lie bracket and its Jacobi identity,
and bilinear actions with the commutator identity and parity preservation.
The purely odd line is constructed as a structure satisfying these laws.
Its dimensions are computed using Mathlib's vector-space
\texttt{Module.finrank}, not introduced as hypotheses.
The integer difference is proved equal to $-1$.
The final theorem \texttt{conjecture\_7400\_false} refutes nonnegativity
even in this restricted class of superrepresentations, which is contained
in the original unrestricted domain.
\end{document}
